\section{An explicit representation of a normal pair}
\label{sec:pairs}
We use the finite colored-graph definition of incidence geometry and the groups
$C(X),K(X)$ from the common paper~\cite{Recognition}. Correlations permute color
classes; $K(X)$ fixes them individually. For a finite group $G$ and a normal
subgroup $H$, the following is Theorem A of D\'iaz Ramos, Molinier and
Viruel~\cite{DMV}, with the rank and connectedness of their construction stated
explicitly.

\begin{theorem}\label{R.FinitePairs}
Every finite normal pair $H\mathrel{\trianglelefteq}G$ is isomorphic to
$(C(\Gamma),K(\Gamma))$ for a finite connected incidence geometry $\Gamma$.
For $G\ne1$ the construction below has rank $[G:H]+3$; for $G=1$ it has rank
one. No fixed rank-two assertion is made.
\end{theorem}

Normality is necessary because $K(\Gamma)$ is the kernel of the palette action.
In particular $G/H$ embeds in a symmetric group on the palette, so rank two
would force $[G:H]\leq2$. This elementary restriction is not evaded by
allowing correlations that exchange types.

\begin{theorem}[Ordered-pair encoding]\label{R.PairEncoding}
Let $G\ne1$ be finite and $H\mathrel{\trianglelefteq}G$. There is a connected
properly colored triangle-free graph $X(G,H)$ of minimum degree two with
\[
 (C(X(G,H)),K(X(G,H)))\cong(G,H).
\]
For $q=|G|$ it has $(q^3+9q^2-8q)/2$ vertices and $[G:H]+3$ colors.
\end{theorem}
\begin{proof}
This is the explicit construction of \cite[Theorem 3.1]{DMV}.
Enumerate $G=\{g_1=1,g_2,\ldots,g_q\}$. Take a vertex $P_x$ for each
$x\in G$. Label the ordered pair $(x,y)$, for $x\ne y$, by $k$ when
$yx^{-1}=g_k$. Thus $k\geq2$. Introduce separate vertices
\[
 T_{x,y},\quad S_{x,y,0},\ldots,S_{x,y,k+2}
\]
for every ordered pair. Suppressing the pair subscripts within one gadget,
add exactly the edges
\begin{equation}\label{eq:arrow}
 P_xS_0,\quad S_0P_y,\quad
 S_iS_{i+1}\ (0\leq i<k+2),\quad
 P_xT,\quad TS_{k+2},\quad S_{k+2}P_y.
\end{equation}
There are no edges between original vertices and no other edges. Color $P_x$
by its coset $xH$. Give all $T$-vertices one new color, the even-indexed
$S$-vertices a second, and the odd-indexed ones a third. All three colors
occur and are disjoint from the coset palette.

Every displayed edge joins different colors: only consecutive chain vertices
use the two parity colors together, and every other edge has distinct vertex
kinds or joins the old and new palettes. A single gadget has no triangle.
Its short paths between the old endpoints have lengths two and three, while
the chain between $S_0,S_{k+2}$ has length $k+2\geq4$. More directly, no
adjacent pair of internal neighbors of either old endpoint occurs in the
edge list, and the internal chain with its one attached $T$ has no triangle.
Distinct gadgets meet only at old vertices, which are pairwise nonadjacent,
so no triangle can use different gadgets. Every pair of old vertices is
connected through its gadget and every new vertex connects to its endpoints.
Thus the graph is connected.

The degrees of old vertices are $4(q-1)$, those of $S_0,S_{k+2}$ are three,
and all remaining vertices have degree two. Consequently the old vertices
are exactly those of degree at least four, including when $q=2$.
Let an uncolored graph automorphism induce $P_x\mapsto P_{\rho(x)}$.
The degree-two neighbors of $P_x$ are precisely the $T_{x,z}$.
Among them $T_{x,y}$ is the unique one at distance two from $P_y$: its other
neighbor is the terminal vertex $S_{x,y,k+2}$, adjacent to exactly that old
vertex $P_y$. The terminal vertex is also the unique degree-three neighbor
of $T_{x,y}$.

Starting at this terminal vertex, the path with degree-two internal vertices
that ends at a degree-three vertex is intrinsic. It is the chain ending at
$S_{x,y,0}$ and has length $k+2$. The other route through $T_{x,y}$ ends at
an old vertex of degree at least four. Hence the graph recovers both the
orientation and the full label $k$ of each ordered pair. Once the old
permutation is fixed, all gadget vertices are forced. Conversely a permutation
of the old vertices preserving these labels extends uniquely by preserving
every gadget index. This is exactly the recognition and extension mechanism
of the common theorem, applied first to the labeled ordered-pair structure.

Label preservation is the equation
\[
 \rho(y)\rho(x)^{-1}=yx^{-1}\qquad(x\ne y).
\]
Taking $x=1$ gives $\rho(y)=y\rho(1)$ for $y\ne1$, and the same equality
holds at $y=1$. Thus the full uncolored automorphism group consists of right
translations. For ordinary composition of permutations use the homomorphism
\[
 g:P_x\longmapsto P_{xg^{-1}},
\]
with the same right translation on both indices of every gadget.
The inverse in this formula fixes the left/right convention; it does not
alter the graph construction.

All translations preserve the three auxiliary colors. Normality of $H$
makes them permute the cosets $xH$, so every uncolored automorphism is a
correlation. A translation fixes all coset colors if and only if $g\in H$:
necessity follows from the coset of the identity, and normality gives
sufficiency on every coset. This proves the exact pair assertion.

For each label $k=2,\ldots,q$ there are $q$ ordered pairs, with $k+4$
internal vertices per pair. The total is
\[
 q+q\sum_{k=2}^{q}(k+4)=\frac{q^3+9q^2-8q}{2}.
\]
This is the size of the borrowed construction, not an optimality statement.
\end{proof}

\begin{proof}[Proof of Theorem~\ref{R.FinitePairs}]
For $G\ne1$, apply the common chamber-completion theorem to
$X(G,H)$. It preserves the correlation pair, palette and connectedness,
so the resulting geometry has rank $[G:H]+3$. For $G=1$, use a single
vertex with one color. Its correlation group and type-preserving subgroup
are both trivial, and its only nonempty clique is a chamber.
\end{proof}

The common cleaning theorem supplies a separate conclusion: any nonempty
finite incidence system, including one with isolated vertices or triangles,
can first be cleaned and then completed without changing its pair. Cleaning
preserves components rather than making a disconnected input connected.
Thus the connectedness asserted above comes from the explicit encoder;
it is not an unstated hypothesis or conclusion of the arbitrary-input case.
